\documentclass[10pt,letterpaper]{article}
\usepackage{technical-implementation}
\setDocTitle{Spacing Poses Easing and Cycles}
\setDocSubtitle{Technical reference | supplementary engineering with book crosswalk}
\setDocVersion{1.0}
\setDocStatus{Educational / Draft}
\setDocOwner{Game Development}
\setDocDate{2026-10-30}
\setDocSummary{Technical reference | supplementary engineering with book crosswalk}
\begin{document}
\makedoctitle
\makedocsummary
\tableofcontents\clearpage

\section{Book foundation and software boundary}
Chapters 27--46 develop spacing, holds, anticipation, follow-through, force cues, and repeat action; Chapters 55--65 develop locomotion, perspective, and motion accents. Formulas below are supplementary mathematical models rather than prescribed drawing methods.
\section{Position, velocity, and spacing}
For sampled positions $x_k$, the displacement $\delta_k=x_{k+1}-x_k$ and approximate speed $|\delta_k|/\Delta t$ change independently of the number of sampled frames. Equal spatial spacing at constant $\Delta t$ produces approximately constant speed. Tight spacing near an extreme with larger spacing mid-action suggests easing.
\section{Normalized interpolation}
Let $u=\operatorname{clamp}((t-t_0)/(t_1-t_0),0,1)$. A linear path uses $E(u)=u$, smoothstep uses $E(u)=3u^2-2u^3$ (zero endpoint velocity), and smootherstep uses $E(u)=6u^5-15u^4+10u^3$ (zero first and second derivatives at endpoints). Evaluate $x(t)=x_0+(x_1-x_0)E(u)$.
\section{Artistic controls}
Anticipation shifts the preparation pose opposite the main intent; follow-through delays dependent parts after main deceleration; overlapping action applies phase offsets to multiple channels. A hold may use stationary global position with subtle expressive motion. Do not apply smoothstep blindly to walk translation, contact states, or a physical ballistic trajectory.
\section{Rotation and wraparound}
A raw scalar interpolation between headings 179 degrees and -179 degrees takes the long way unless angles are wrapped or represented by quaternions. For $S^1$, interpolate the signed shortest angular difference. For $SO(3)$ use normalized quaternions and spherical interpolation with a consistent sign convention.
\section{Cycles and contact integrity}
For a periodic clip of duration $P>0$, define $u=\operatorname{fract}(t/P)$. Check continuity at the seam for pose \emph{and} velocity where intended. In world-space movement a contact foot must not slide: compare its world transform before, during, and after contact rather than measuring only its position relative to the character root.
\section{Example measurement}
\[e_{\mathrm{contact}}=\max_{k\in C}\|p_{\mathrm{foot}}(t_k)-p_{\mathrm{anchor}}\|,\qquad e_{\mathrm{seam}}=\|p(0)-p(P^-)\|.\]
Set tolerances in project units, measured at the target runtime step. A zero seam position error alone does not guarantee a smooth-loop velocity transition.
\section{Review protocol}
View at 100\% speed, quarter-speed, and as a contact sheet. Compare alternative ease profiles while keeping shot length and key extremes fixed. Annotate at least one intentional violation of physical plausibility and the experience goal it supports.
\end{document}
